\chapter*{Περίληψη}
\addcontentsline{toc}{chapter}{Περίληψη}

% TODO: Η περίληψη είναι ζωντανό κείμενο - ενημερώστε την όταν ολοκληρωθούν
% τα επόμενα βήματα (βλ. Κεφάλαιο 6, Μελλοντική Εργασία).

Η σύγκλιση επίγειων δικτύων 5G/6G (Terrestrial Networks - TN) με δίκτυα
δορυφόρων χαμηλής τροχιάς (Low Earth Orbit - LEO) σε Ολοκληρωμένα
Δορυφορικά-Επίγεια Δίκτυα (Satellite-Terrestrial Integrated Networks -
STINs) αποτελεί βασικό σχεδιαστικό στόχο του 6G, με στόχο την παροχή
καθολικής κάλυψης πέρα από τους περιορισμούς της επίγειας υποδομής. Η
πολυσυνδεσιμότητα (Multi-Connectivity - MC) μεταξύ επίγειων σταθμών βάσης
και δορυφόρων προτείνεται ως μηχανισμός βελτίωσης της αξιοπιστίας, της
κάλυψης και της διαχείρισης φορτίου σε τέτοια δίκτυα.

Η παρούσα διπλωματική εργασία αναπτύσσει ένα πλαίσιο προσομοίωσης σε
MATLAB για ένα σενάριο συνδεσιμότητας δύο επίγειων σταθμών βάσης (5G NR,
3GPP TR 38.901 UMa/UMi) και ενός δορυφόρου LEO, στο οποίο κάθε χρήστης
επιλέγει τον κόμβο εξυπηρέτησης (επίγειο ή δορυφορικό) με βάση το
υψηλότερο SNR (greedy single-connectivity επιλογή). Το μοντέλο υπολογίζει
απώλειες διάδοσης βάσει 3GPP TR 38.901 (με στοχαστική επιλογή LOS/NLOS
ανά ζεύξη και log-normal shadow fading) για το επίγειο σκέλος, και
ελεύθερου χώρου (FSPL) με μάσκα ελάχιστης γωνίας ανύψωσης για το
δορυφορικό σκέλος, και εξάγει δείκτες απόδοσης (KPIs) ανά χρήστη:
χωρητικότητα Shannon, και ενέργεια ανά bit βάσει γραμμικού μοντέλου
ισχύος (EARTH για τον σταθμό βάσης, μοντέλο απόδοσης ενισχυτή ισχύος για
τον δορυφόρο).

Πέρα από το στατικό σενάριο αναφοράς, η εργασία επεκτείνει την
προσομοίωση σε δύο κατευθύνσεις: (α) επαναλαμβανόμενες εκτελέσεις της
ίδιας τοπολογίας με διαφορετική στοχαστική πραγματοποίηση καναλιού, ώστε
να χαρακτηριστεί η διακύμανση των KPI, και (β) χρονική εξέλιξη πολλαπλών
διαδοχικών μεταδόσεων κατά τη διάρκεια ενός πάσσου του δορυφόρου LEO,
όπου το υποδορυφορικό σημείο μετακινείται με πραγματική ταχύτητα ίχνους
εδάφους υπολογισμένη από την Κεπλεριανή τροχιακή περίοδο, ώστε να
αποτυπωθεί η μεταβολή της γωνίας ανύψωσης, του καναλιού, και των
handovers μεταξύ κόμβων με τον χρόνο. Επιπλέον, ένα πρώτο πλαίσιο Monte
Carlo παράγει ένα ετικετοποιημένο σύνολο δεδομένων πάνω σε τυχαιοποιημένες
τοπολογίες, το οποίο χρησιμοποιείται για μια πρώτη δοκιμή μοντέλων
μηχανικής μάθησης (Λογιστική Παλινδρόμηση, Random Forest) που προβλέπουν
τον τύπο εξυπηρέτησης (επίγειος/δορυφορικός) από τα χαρακτηριστικά των
υποψήφιων ζεύξεων.

Τα αποτελέσματα δείχνουν σαφή διαφοροποίηση μεταξύ επίγειας και
δορυφορικής εξυπηρέτησης ως προς τα τρία εξεταζόμενα KPI (throughput,
ενέργεια ανά bit, SNR), με τον δορυφόρο να προσφέρει χαμηλότερο ρυθμό
αλλά σημαντικά καλύτερη ενεργειακή απόδοση ανά bit. Η χρονική
προσομοίωση πάσσου αναδεικνύει δύο περιορισμούς του τρέχοντος μοντέλου
επιλογής κόμβου που δεν ήταν ορατοί στο στατικό σενάριο: (i) τεχνητά
συχνά handovers λόγω μη-συσχετισμένου (i.i.d.) shadow fading ανά χρονικό
βήμα, και (ii) απουσία κατάστασης "outage" όταν κανένας υποψήφιος κόμβος
δεν προσφέρει χρησιμοποιήσιμο SNR. Τα ευρήματα αυτά καθορίζουν τη
μελλοντική κατεύθυνση της εργασίας: χρονικά συσχετισμένο μοντέλο shadow
fading, ρητή διαχείριση outage, και επέκταση του μοντέλου μηχανικής
μάθησης πέρα από την τετριμμένη πρόβλεψη argmax(SNR) προς πιο ρεαλιστικούς
στόχους πρόβλεψης.

\vspace{1em}
\noindent\textbf{Λέξεις Κλειδιά:} πολυσυνδεσιμότητα, ολοκληρωμένα
δορυφορικά-επίγεια δίκτυα, 6G, δορυφόροι LEO, 3GPP TR 38.901, μηχανική
μάθηση, ενεργειακή απόδοση, handover
